% Teaching environments.
%
% In the reading flow: thought experiments (and their return at the end of a
% topic), worked examples, key ideas and checks. Beside the reading flow, as
% floating sidebars: history and going deeper.

\tcbset{
  sebox/.style={
    enhanced jigsaw, breakable, sharp corners, frame hidden, boxrule=0pt,
    left=2.6mm, right=2.6mm, top=1.0mm, bottom=1.8mm, boxsep=0pt,
    lefttitle=2.6mm, righttitle=2.6mm, toptitle=1.7mm, bottomtitle=0.3mm,
    before skip=1.0\baselineskip plus 3pt minus 2pt,
    after skip=1.0\baselineskip plus 3pt minus 2pt,
    fonttitle=\sffamily\bfseries\small, titlerule=0pt,
    parbox=false, before upper={\setlength{\parindent}{1em}},
  },
  sebar/.style={
    sebox, colback=#1!9, colbacktitle=#1!9, coltitle=#1!82!black,
    borderline west={2.4pt}{0pt}{#1},
  },
  sesidebar/.style={
    enhanced, sharp corners, frame hidden, boxrule=0pt,
    left=2.6mm, right=2.6mm, top=1.0mm, bottom=1.8mm, boxsep=0pt,
    lefttitle=2.6mm, righttitle=2.6mm, toptitle=1.7mm, bottomtitle=0.3mm,
    fonttitle=\sffamily\bfseries\small, titlerule=0pt, fontupper=\small,
    parbox=false, before upper={\setlength{\parindent}{1em}},
    colback=#1!8, colbacktitle=#1!8, coltitle=#1!85!black,
    borderline north={1.2pt}{0pt}{#1},
    float=tb,
  },
}

% A puzzle that opens a topic. The topic gives the reader what is needed to
% solve it, and the topic closes by returning to it. Short boxes (thought
% experiments and their return, key ideas, checks) never break, so a topic
% heading moves to the next column together with its thought experiment.
\makeatletter
\newcommand{\seafterbox}{\par\@afterindentfalse\@afterheading}
\makeatother
\newtcolorbox{thought}[1]{
  sebar=seorange, unbreakable, after app={\seafterbox},
  title={\faLightbulb[regular]\enspace Thought experiment\enspace
    {\color{black!85}\textperiodcentered\enspace #1}},
}
\newtcolorbox{revisited}[1]{
  sebox, unbreakable, colback=white, colbacktitle=white, coltitle=seorange!82!black,
  borderline west={2.4pt}{0pt}{seorange},
  title={\faLightbulb\enspace Back to the thought experiment\enspace
    {\color{black!85}\textperiodcentered\enspace #1}},
}

% Worked examples, numbered within chapters.
\newtcolorbox[auto counter, number within=chapter,
  crefname={Example}{Examples}, Crefname={Example}{Examples}]{example}[2][]{
  sebox, colback=seblue!5, colbacktitle=seblue!5, coltitle=seblue,
  borderline north={1.3pt}{0pt}{seblue},
  title={Example~\thetcbcounter\hspace{0.7em}{\color{black!85}#2}},
  #1,
}
\newcommand{\solution}{\par\smallskip\noindent
  {\sffamily\bfseries\small\color{seblue}Solution}\enspace}
\newcommand{\discussion}{\par\smallskip\noindent
  {\sffamily\bfseries\small\color{seblue}What it tells us}\enspace}

% A central idea stated once, plainly.
\newtcolorbox{keyidea}[1][]{
  sebar=seblue, unbreakable,
  title={\faKey\enspace Key idea\ifstrempty{#1}{}{\enspace{\color{black!85}\textperiodcentered\enspace #1}}},
}

% Short questions in the text; answers at the end of each chapter.
\newtcolorbox[auto counter, number within=chapter,
  crefname={Check}{Checks}, Crefname={Check}{Checks}]{checkpoint}[1][]{
  sebar=seteal, unbreakable,
  title={\faCheckCircle[regular]\enspace Check your understanding~\thetcbcounter},
  #1,
}

% Sidebars beside the reading flow.
\newtcolorbox{deeper}[1]{
  sesidebar=seviolet,
  title={Going deeper\enspace{\color{black!85}\textperiodcentered\enspace #1}},
}
\newtcolorbox{history}[1]{
  sesidebar=sesand,
  title={History\enspace{\color{black!85}\textperiodcentered\enspace #1}},
}

% ---------------------------------------------------------------------------
% Chapter openings: a lead paragraph beside the list of goals, across both
% columns. Use before \chapter:
%   \chapteropening{lead text}{\item goal \item goal ...}
\newtcolorbox{goalsbox}{
  enhanced, sharp corners, frame hidden, boxrule=0pt,
  colback=seblue!6, colbacktitle=seblue!6, coltitle=seblue, titlerule=0pt,
  left=3mm, right=3mm, top=0.6mm, bottom=2mm, toptitle=2mm, bottomtitle=0.2mm,
  lefttitle=3mm, fonttitle=\sffamily\bfseries\small, fontupper=\small,
  borderline north={1.3pt}{0pt}{seblue},
  title={In this chapter you will learn},
}
\newcommand{\chapteropening}[2]{%
  \setchapterpreamble[u]{%
    \vspace*{0.1\baselineskip}%
    \noindent
    \begin{minipage}[t]{0.47\textwidth}\vspace{0pt}%
      \setlength{\parindent}{0pt}\setlength{\parskip}{0.6\baselineskip}%
      {\fontsize{11.3}{15.2}\selectfont #1\par}%
    \end{minipage}\hfill
    \begin{minipage}[t]{0.49\textwidth}\vspace{0pt}%
      \begin{goalsbox}
        \begin{itemize}[leftmargin=1.1em, itemsep=0.35ex, topsep=0.2ex]
          #2
        \end{itemize}
      \end{goalsbox}
    \end{minipage}\par
    \vspace{1.3\baselineskip}%
  }%
}

% ---------------------------------------------------------------------------
% End-of-chapter sections.
\newcommand{\closingsection}[1]{\section*{#1}\markright{#1}}
\newenvironment{chaptersummary}
  {\closingsection{Summary}\begin{itemize}}
  {\end{itemize}}
\newenvironment{checkanswers}
  {\closingsection{Answers to Check your understanding}%
   \begin{description}[style=nextline, font=\normalfont\sffamily\bfseries\small\color{seteal}]}
  {\end{description}}
\newlist{problemlist}{enumerate}{1}
\setlist[problemlist]{
  label={\sffamily\bfseries\small\color{seblue}\thechapter.\arabic*},
  ref=\thechapter.\arabic*, leftmargin=*, labelsep=0.6em, itemsep=0.8ex,
}
\crefname{problemlisti}{Problem}{Problems}
\Crefname{problemlisti}{Problem}{Problems}
\newenvironment{problems}
  {\closingsection{Problems}\begin{problemlist}}
  {\end{problemlist}}
\newenvironment{furtherreading}
  {\closingsection{Sources and further reading}%
   \begin{description}[style=nextline, font=\normalfont]}
  {\end{description}}
